% !PRINTED 691
\weberpage{691}

%<<CONT
is a function in $\Omega$, which we call a \gesperrt{differential quotient}
of the first kind.

If $\mathfrak{P}$ is a point at which $(z - z_0)$ becomes infinitely small of
order $e$, then this point occurs $(e - 1)$ times in $\mathfrak{Z}$,
and consequently $[w(z - z_0)]_0 = 0$. If $z$ becomes infinite in
another point $\mathfrak{P}$ in the $e$th order, then this point occurs
$e$ times in $\mathfrak{U}$ and $(e - 1)$ times in $\mathfrak{Z}$, hence once in the
numerator of $wz$, and consequently at this point too $(wz)_0 = 0$,
and these two requirements, namely:

a) At every point at which $z$ has a finite value $z_0$,
$[w(z - z_0)]_0 = 0$;

b) At a point at which $z$ becomes infinite, $(wz)_0 = 0$,
are also sufficient to define a differential quotient of the first kind
$w$.

In order to satisfy the first condition, we take a
minimal basis with respect to $z$:
\weq{2}{\omega_1,\ \omega_2,\ \dots,\ \omega_n,}
put
\weq{3}{a_{r,\,s} = S(\omega_r\omega_s),}
so that
\weq{4}{\Delta = \Sigma \pm a_{1,1}a_{2,2}\dots a_{n,n}}
is the field discriminant (with respect to $z$), and determine a basis
\weq{5}{\varepsilon_1,\ \varepsilon_2,\ \dots,\ \varepsilon_n}
from the equations:
\weq{6}{\omega_r = \Sigma a_{r,\,i}\varepsilon_i.}

If $(r, s)$ is a symbol which is $= 0$ when $r$ and $s$ are different,
and $= 1$ when $r = s$, then one can determine the rational
functions $a'_{r,\,s}$ of $z$ in such a way that
\weq{7}{\sideset{^i}{}\sum a_{i,\,r}a'_{i,\,s} = (r,\ s)}
holds, and one obtains from (6):
\weq{8}{\varepsilon_r = \sideset{^i}{}\sum a'_{i,\,r}\omega_i,}
and the quantities $a'_{r,\,s}$ have no denominator other than $\Delta$.
From this there follows, by reason of (3) and (7):
\weq{9}{\begin{aligned}
S(\varepsilon_r\omega_s) &= (r,\ s), \\
a'_{r,\,s} &= S(\varepsilon_r\varepsilon_s).
\end{aligned}}
